\section{Alineamiento ciclotómico, orbifold de orden \(3\) y Moonshine monstruoso}
\label{p12-83-c20-s8:chap:alineamiento-orbifold-moonshine}

El retículo de Leech construido en la
\cref{p12-83-c20-s5:thm:identificacion-vecino-leech} permite pasar de la incidencia
finita a una estructura graduada infinita. Esta transición no se efectuará
mediante una cadena de nombres. Se construirán, en orden tipado:
\begin{enumerate}
 \item el entrelazamiento de las doce fibras \(A_2\);
 \item una isometría reticular fijo-libre de orden tres;
 \item su traza graduada y su elevación de tipo cero;
 \item el orbifold cíclico y su automorfismo dual de clase \(3B\);
 \item el álgebra de operadores de vértice lunar;
 \item por publicaciones distintas, su grupo de automorfismos, su carácter
       graduado y las series de McKay--Thompson.
\end{enumerate}
Moonshine será el teorema que coordina la acción y las series. No se usará
el atajo no tipado
\(V^\natural\to\mathbb M\to\text{Moonshine}\).

\subsection{Alineamiento de las \(12\) fibras APP y Witt}

Las doce fibras de \(W_{\mathrm{APP}}\cong A_2^{12}\) llevan tres
etiquetas independientes:
\[
 i\in\mathbb Z/3\mathbb Z
 \quad\text{(par de coordenadas)},\qquad
 \mu\in\{0,1\}
 \quad\text{(hoja)},\qquad
 \varepsilon\in\{0,1\}
 \quad\text{(órbita ciclotómica)}.
\]
Fijado el origen dodecafásico construido en el TPK, definimos
\begin{equation}
 \lambda(i,\mu,\varepsilon)
 :=4i+2\mu+\varepsilon\pmod{12}.
 \label{p12-83-c20-s8:eq:alineamiento-doce-fibras-u031}
\end{equation}

\begin{lemma}[Biyectividad del alineamiento]
\label{p12-83-c20-s8:lem:biyectividad-alineamiento-u031}
La aplicación \(\lambda\) de
\eqref{p12-83-c20-s8:eq:alineamiento-doce-fibras-u031} es una biyección sobre
\(\mathbb Z/12\mathbb Z\).
\end{lemma}

\begin{proof}
Para cada \(i\), los cuatro valores \(4i+2\mu+\varepsilon\) forman el bloque
\(\{4i,4i+1,4i+2,4i+3\}\). Los tres bloques para \(i=0,1,2\) son disjuntos
y cubren las doce clases.
\end{proof}

Sea \(C\) el Coxeter de \(A_2\) y sea \(\mathsf R\) la reflexión que
satisface
\begin{equation}
 \mathsf R C\mathsf R=C^{-1}.
 \label{p12-83-c20-s8:eq:reflexion-conjuga-coxeter-u031}
\end{equation}
Fijadas las trivializaciones \(\iota_b\) y los marcos orientados \(P_b\),
ponemos
\begin{align}
 g_b&:=\iota_b^{-1}P_bCP_b^{-1}\iota_b,
 \label{p12-83-c20-s8:eq:coxeter-calibrado-fibra-u031}\\
 T_{i\mu\varepsilon}
 &:=\iota_{\lambda(i,\mu,\varepsilon)}^{-1}
 P_{\lambda(i,\mu,\varepsilon)}\mathsf R^\mu.
 \label{p12-83-c20-s8:eq:entrelazador-fibra-u031}
\end{align}

\begin{theorem}[Entrelazamiento APP--Witt]
\label{p12-83-c20-s8:thm:alineamiento-app-witt-u031}
Para cada terna \((i,\mu,\varepsilon)\),
\begin{equation}
 g_{\lambda(i,\mu,\varepsilon)}T_{i\mu\varepsilon}
 =T_{i\mu\varepsilon}C^{(-1)^\mu}.
 \label{p12-83-c20-s8:eq:identidad-entrelazamiento-app-witt-u031}
\end{equation}
La suma directa define un isomorfismo \(C_3\)-equivariante
\begin{equation}
 T_\lambda:
 W_{\mathrm{APP}}\otimes\mathbb C
 \xrightarrow{\;\cong\;}
 W_c:=\bigoplus_{b=0}^{11}
 (A_{2,b}\otimes\mathbb C).
 \label{p12-83-c20-s8:eq:isomorfismo-app-witt-u031}
\end{equation}
\end{theorem}

\begin{proof}
Para \(\mu=0\), la definición de \(g_b\) da directamente la conjugación
por \(C\). Para \(\mu=1\), se inserta
\(\mathsf RC\mathsf R=C^{-1}\). El
\cref{p12-83-c20-s8:lem:biyectividad-alineamiento-u031} garantiza que la suma directa
usa exactamente una vez cada posición Witt y conserva par, hoja y órbita.
\end{proof}

El teorema es relativo al transporte simultáneo del orden de pares, hojas y
órbitas, del origen Witt, de los marcos y de la orientación. Cambiar una
sola de estas coordenadas no define el mismo funtor.

\subsection{Traza graduada del generador ternario}

Sea
\[
 \rho_W:=\bigoplus_{b=0}^{11}g_b
\]
y consideremos la envolvente bosónica
\begin{equation}
 \mathcal H_{\rm bos}(W_c)
 :=
 \operatorname{Sym}\!\left(
 \bigoplus_{n\ge1}W_c[n]
 \right).
 \label{p12-83-c20-s8:eq:fock-bosonico-doce-fibras-u031}
\end{equation}
Si \(N\) es el operador de grado, definimos
\begin{equation}
 Z_{\rm Fock}(q)
 :=
 q^{-1}\operatorname{Tr}_{\mathcal H_{\rm bos}(W_c)}
 \left(\Gamma_{\rm bos}(\rho_W)q^N\right).
 \label{p12-83-c20-s8:eq:traza-graduada-fock-u031}
\end{equation}

\begin{proposition}[Producto ciclotómico de nivel tres]
\label{p12-83-c20-s8:prop:producto-tres-doce-fibras-u031}
Como serie formal de Laurent,
\begin{equation}
 \boxed{
 Z_{\rm Fock}(q)
 =q^{-1}\prod_{n\ge1}(1+q^n+q^{2n})^{-12}
 =:t_3(q).}
 \label{p12-83-c20-s8:eq:t3-producto-u031}
\end{equation}
Su expansión comienza por
\begin{equation}
 t_3(q)
 =q^{-1}-12+54q-76q^2-243q^3+1188q^4+O(q^5).
 \label{p12-83-c20-s8:eq:t3-expansion-u031}
\end{equation}
\end{proposition}

\begin{proof}
Cada \(g_b\) es conjugado a \(C\), por lo que
\[
 \det(I-g_bq^n)=1+q^n+q^{2n}.
\]
La identidad
\[
 \sum_{k\ge0}
 \operatorname{Tr}(\operatorname{Sym}^k\rho_W)t^k
 =\det(I-\rho_Wt)^{-1}
\]
aplicada en cada grado produce el producto. La multiplicación formal hasta
grado cinco da la expansión. En particular, el coeficiente de \(q\) se
obtiene como
\[
 [q^2]\,
 (1+q+q^2)^{-12}(1+q^2+q^4)^{-12}
 =66-12=54.
\]
\end{proof}

El valor \(54\) aparece aquí como coeficiente de una traza construida desde
doce fibras. La unidad U030 lo obtuvo también como índice orientado del
agregado APP y probó que doce es la única solución positiva de la rigidez
ciclotómica--radial. La igualdad de ambos valores es una compatibilidad de
dos realizaciones; ninguna usa el coeficiente de \(J\) como entrada.

\subsection{Palabra de soporte total e isometría fijo-libre}

Sea
\[
 q_*=(1,1,1,1,1,1)\in\mathbb F_3^6.
\]
La multiplicación por la matriz de Paley--Witt da
\begin{equation}
 q_*A_W=(2,1,1,1,1,1).
 \label{p12-83-c20-s8:eq:producto-qstar-aw-u031}
\end{equation}
Por tanto,
\begin{equation}
 c_*:=(q_*,q_*A_W)
 =(1,1,1,1,1,1,2,1,1,1,1,1)
 \in C_W
 \label{p12-83-c20-s8:eq:palabra-soporte-total-u031}
\end{equation}
y todas sus coordenadas son no nulas.

Para cada posición \(b\), fijamos
\begin{equation}
 P_b(c_*)=
 \begin{cases}
  \mathsf R,&(c_*)_b=1,\\
  I,&(c_*)_b=2,
 \end{cases}
 \qquad
 g_c:=\bigoplus_{b=0}^{11}P_b(c_*)CP_b(c_*)^{-1}.
 \label{p12-83-c20-s8:eq:isometria-orden-tres-u031}
\end{equation}

\begin{proposition}[Orden, ausencia de puntos fijos y pegado]
\label{p12-83-c20-s8:prop:isometria-fijo-libre-pegado-u031}
El operador \(g_c\) tiene orden tres, carece de vectores fijos no nulos y
preserva el retículo de Niemeier \(N(A_2^{12})\) construido en U021.
\end{proposition}

\begin{proof}
Cada bloque es conjugado al Coxeter \(C\), cuyo polinomio mínimo es
\(X^2+X+1\). Luego tiene orden tres y no posee autovalor uno. La suma
directa conserva ambas propiedades. El Coxeter actúa trivialmente sobre
\(A_2^*/A_2\); por ello cada bloque conserva la clase discriminante y la
unión de clases laterales que define el pegado.
\end{proof}

\subsection{Preservación del vecino de Leech}

Escribamos
\[
 v:=v_{\mathcal F_*}
 =4\rho_1+\rho_2+\cdots+\rho_{12},
 \qquad
 v^2=54,
\]
para el vector seleccionado por el origen incidencial de
\eqref{p12-83-c20-s5:eq:vector-marco-incidencial}. Sean
\[
 N:=N(A_2^{12}),
 \qquad
 N_0:=\ker(\langle\,\cdot\,,v\rangle\bmod3),
 \qquad
 N_v:=N_0+\mathbb Z\frac v3\cong\Lambda_{24}.
\]

\begin{theorem}[Preservación orientada del vecino]
\label{p12-83-c20-s8:thm:preservacion-vecino-orden-tres-u031}
Los desplazamientos
\begin{equation}
 y_*:=\frac{g_cv-v}{3},
 \qquad
 z_*:=\frac{g_c^{-1}v-v}{3}
 \label{p12-83-c20-s8:eq:desplazamientos-vecino-u031}
\end{equation}
pertenecen a \(N_0\). En consecuencia,
\[
 g_cN_0=N_0,
 \qquad
 g_cN_v=N_v.
\]
\end{theorem}

\begin{proof}
Las palabras discriminantes de \(y_*\) y \(z_*\) son \(c_*\) y
\(-c_*\), ambas en \(C_W\); de ahí \(y_*,z_*\in N\). El cálculo en cada
factor \(A_2\) da
\begin{equation}
 \langle y_*,v\rangle
 =\langle z_*,v\rangle
 =-(4^2+11)=-27\equiv0\pmod3.
 \label{p12-83-c20-s8:eq:emparejamientos-desplazamientos-u031}
\end{equation}
Luego ambos están en \(N_0\). Para \(x\in N_0\),
\[
 \langle g_cx,v\rangle
 =\langle x,g_c^{-1}v\rangle
 =\langle x,v\rangle+3\langle x,z_*\rangle
 \equiv0\pmod3.
\]
La misma cuenta con \(g_c^{-1}\) da igualdad del núcleo. Finalmente,
\[
 g_c(v/3)=v/3+y_*,
\]
por lo que se preserva la clase que genera el vecino.
\end{proof}

\subsection{\(24\) orientaciones y naturalidad reticular}

El enumerador de pesos de \eqref{p12-83-c20-s1:eq:enumerador-pesos-cw} contiene
\(24y^{12}\). Por tanto, el conjunto
\[
 C_W^\times:=
 \{u\in C_W:u_b\in\{1,2\}\text{ para todo }b\}
\]
posee 24 elementos. Para cada \(u\in C_W^\times\), definimos
\[
 P_b(u)=
 \begin{cases}
  \mathsf R,&u_b=1,\\
  I,&u_b=2,
 \end{cases}
 \qquad
 g_u:=\bigoplus_{b=0}^{11}P_b(u)CP_b(u)^{-1}.
\]

\begin{proposition}[Robustez de las orientaciones de soporte total]
\label{p12-83-c20-s8:prop:robustez-orientaciones-u031}
Para toda \(u\in C_W^\times\),
\[
 \frac{g_uv-v}{3},
 \quad
 \frac{g_u^{-1}v-v}{3}
 \in N_0.
\]
En particular, cada \(g_u\) preserva el vecino construido con la
orientación transportada.
\end{proposition}

\begin{proof}
En coordenadas de raíces simples,
\[
 C\rho-\rho=(-2,-1),
 \qquad
 C^{-1}\rho-\rho=(-1,-2).
\]
Las clases discriminantes de los dos desplazamientos son \(u\) y \(-u\).
Para el coeficiente radial \(a_b\in\{4,1\}\),
\[
 \left\langle
 \frac{(C^{\pm1}-I)a_b\rho}{3},a_b\rho
 \right\rangle=-a_b^2.
\]
La suma sobre los doce factores vuelve a dar \(-27\). La prueba del
\cref{p12-83-c20-s8:thm:preservacion-vecino-orden-tres-u031} se aplica componente a
componente.
\end{proof}

Para tipar la recalibración, sea
\(h\in\operatorname{Aut}(C_W)\) monomial: una permutación \(\sigma\) de
las doce coordenadas seguida de multiplicadores
\(\epsilon_b\in\{1,2\}\). Definimos su elevación reticular
\begin{equation}
 \widetilde h
 :=P_\sigma\circ
 \bigoplus_{b=0}^{11}\mathsf R^{\nu_b},
 \qquad
 \nu_b=
 \begin{cases}
 0,&\epsilon_b=1,\\
 1,&\epsilon_b=2.
 \end{cases}
 \label{p12-83-c20-s8:eq:elevacion-monomial-u031}
\end{equation}
Su acción en el discriminante es exactamente \(h\). Si el cambio transporta
simultáneamente código, bandera, origen incidencial, vector \(v\), vecino y
marcos, entonces
\begin{equation}
 g_{hu}=\widetilde h\,g_u\,\widetilde h^{-1}.
 \label{p12-83-c20-s8:eq:naturalidad-isometrias-u031}
\end{equation}
La matriz ternaria \(h\) no actúa sin elevación sobre \(A_2^{12}\) ni sobre
Leech.

\subsection{Traza reticular y cálculo del tipo \(0\)}

Sea \(\widehat g_c\) la elevación estándar de \(g_c\) al VOA reticular
\[
 V_{N_v}=M(1)\otimes\mathbb C_\varepsilon[N_v].
\]
Para \(j=1,2\), la descomposición oscilador--retículo produce
\begin{equation}
 \begin{split}
 Z_{0,j}(\tau)
 &:=
 \operatorname{Tr}_{V_{N_v}}
 \left(\widehat g_c^{\,j}q^{L_0-1}\right)\\
 &=q^{-1}
 \left(
 \sum_{\lambda\in N_v^{g_c^j}}
 \varepsilon_j(\lambda)
 q^{\langle\lambda,\lambda\rangle/2}
 \right)
 \prod_{n\ge1}\det(I-g_c^jq^n)^{-1}.
 \end{split}
 \label{p12-83-c20-s8:eq:factorizacion-traza-reticular-u031}
\end{equation}

\begin{proposition}[Traza de los sectores no triviales]
\label{p12-83-c20-s8:prop:traza-reticular-t3-u031}
Se cumple
\[
 Z_{0,1}(\tau)=Z_{0,2}(\tau)=t_3(\tau).
\]
\end{proposition}

\begin{proof}
La ausencia de puntos fijos implica
\[
 N_v^{g_c}=N_v^{g_c^2}=\{0\},
\]
de modo que la suma reticular de
\eqref{p12-83-c20-s8:eq:factorizacion-traza-reticular-u031} vale uno. En cada factor,
\(g_c^j\) tiene autovalores \(\zeta_3,\zeta_3^2\); por tanto,
\[
 \det(I-g_c^jq^n)=1+q^n+q^{2n}.
\]
Los doce factores reproducen \eqref{p12-83-c20-s8:eq:t3-producto-u031}.
\end{proof}

\begin{theorem}[Peso torcido y tipo de la elevación]
\label{p12-83-c20-s8:thm:peso-torcido-tipo-cero-u031}
La elevación estándar tiene orden tres, peso torcido
\begin{equation}
 \rho_{g_c}=\frac43
 \label{p12-83-c20-s8:eq:peso-torcido-u031}
\end{equation}
y tipo cero.
\end{theorem}

\begin{proof}
En
\(\mathfrak h=N_v\otimes_{\mathbb Z}\mathbb C\), cada uno de los doce
factores aporta un autovector de autovalor \(\zeta_3\) y otro de autovalor
\(\zeta_3^2\). Por tanto,
\[
 \dim\mathfrak h_{\zeta_3}
 =\dim\mathfrak h_{\zeta_3^2}=12.
\]
La fórmula del peso del vacío torcido de orden tres da
\[
 \rho_{g_c}
 =\frac1{4\cdot3^2}
 \bigl(1\cdot2\cdot12+2\cdot1\cdot12\bigr)
 =\frac43.
\]
Para una isometría reticular de orden impar, la elevación estándar conserva
el orden. El tipo es la clase de
\[
 3^2\rho_{g_c}=12\equiv0\pmod3.
\]
Así, las dos hipótesis inicialmente independientes quedan calculadas aquí.
\end{proof}

\subsection{Orbifold cíclico y clase dual \texorpdfstring{\(3B\)}{3B}}

\begin{theorem}[Orbifold triádico del vecino de Leech]
\label{p12-83-c20-s8:thm:orbifold-triadico-vnatural-u031}
Los teoremas clásicos de orbifold cíclico, aplicados a la elevación estándar
del \cref{p12-83-c20-s8:thm:peso-torcido-tipo-cero-u031}, dan
\begin{equation}
 \operatorname{Orb}_{\mathbb Z/3}
 (V_{N_v},\widehat g_c)
 \cong V^\natural.
 \label{p12-83-c20-s8:eq:orbifold-triadico-vnatural-u031}
\end{equation}
El automorfismo dual pertenece a la clase \(3B\) y su serie es
\begin{equation}
 T_{3B}(\tau)=t_3(\tau)+12.
 \label{p12-83-c20-s8:eq:serie-3b-u031}
\end{equation}
\end{theorem}

\begin{proof}
La isometría es fijo-libre de orden tres; el peso torcido y el tipo cero se
calcularon, no se supusieron. Los teoremas de orbifold cíclico identifican
la extensión holomorfa con el módulo lunar y clasifican el automorfismo dual
como \(3B\)
\cite{ChenLamShimakura2016,AbeLamYamada2017,Carnahan2017}. La
identificación usa estos teoremas después de construir el objeto reticular;
no se deduce sólo de una igualdad de series.
\end{proof}

Si \(V^{(i)}\) es el módulo \(\widehat g_c^{\,i}\)-torcido y
\[
 Z_{i,j}(\tau):=
 \operatorname{Tr}_{V^{(i)}}\!
 \left(
 \phi_i(\widehat g_c)^j q^{L_0-1}
 \right),
\]
el carácter del orbifold es el promedio de los nueve sectores:
\begin{equation}
 \operatorname{ch}_{V^\natural}(\tau)
 =\frac13\sum_{i,j\in\mathbb Z/3}Z_{i,j}(\tau)
 =J(\tau).
 \label{p12-83-c20-s8:eq:promedio-nueve-sectores-u031}
\end{equation}
La identidad racional de nivel tres es
\begin{equation}
 J
 =t_3+12+\frac{10\,3^9}{t_3}
 +\frac{4\,3^{14}}{t_3^2}
 +\frac{3^{18}}{t_3^3}.
 \label{p12-83-c20-s8:eq:j-desde-t3-u031}
\end{equation}
Como \(t_3^{-1}=q+O(q^2)\),
\begin{equation}
 [q]J=54+10\cdot3^9=196884
 =54(5\cdot729+1).
 \label{p12-83-c20-s8:eq:coeficiente-196884-u031}
\end{equation}
Esta es una compatibilidad graduada dentro del corredor construido, no un
morfismo escalar \(54\to V^\natural\).

\subsection{Construcción reticular de orden \(2\)}

La ruta de orden tres anterior y la construcción clásica de
Frenkel--Lepowsky--Meurman producen dos presentaciones del mismo VOA lunar.
Para declarar la segunda, fijamos la extensión central
\begin{equation}
 1\longrightarrow\{\pm1\}
 \longrightarrow\widehat\Lambda_{24}
 \longrightarrow\Lambda_{24}
 \longrightarrow1
 \label{p12-83-c20-s8:eq:extension-central-leech-u031}
\end{equation}
con conmutador
\((-1)^{\langle\lambda,\mu\rangle}\), y un cociclo
\(\varepsilon(\lambda,\mu)\) que la representa. El álgebra de grupo
torcida es \(\mathbb C_\varepsilon[\Lambda_{24}]\), y
\begin{equation}
 V_\Lambda
 :=M(1)\otimes\mathbb C_\varepsilon[\Lambda_{24}]
 \label{p12-83-c20-s8:eq:voa-reticular-leech-u031}
\end{equation}
es el VOA reticular de carga central 24.

La negación \(\lambda\mapsto-\lambda\) posee una elevación
\(\theta\) compatible con \eqref{p12-83-c20-s8:eq:extension-central-leech-u031}. Su traza
es
\begin{equation}
 t_2(\tau)
 :=
 \operatorname{Tr}_{V_\Lambda}
 (\theta q^{L_0-1})
 =q^{-1}\prod_{n\ge1}(1+q^n)^{-24}.
 \label{p12-83-c20-s8:eq:t2-leech-u031}
\end{equation}
En efecto, la negación empareja \(e^\lambda\) con \(e^{-\lambda}\), y la
ausencia de torsión deja sólo \(\lambda=0\) en la traza reticular; cada uno
de los 24 osciladores aporta \((1+q^n)^{-1}\).

Sea \(V_\Lambda^T\) el módulo irreducible torcido por \(\theta\), y
\((V_\Lambda^T)^+\) su sector de signo positivo. La construcción FLM es
\begin{equation}
 V^\natural
 =(V_\Lambda)^+\oplus(V_\Lambda^T)^+.
 \label{p12-83-c20-s8:eq:vnatural-flm-u031}
\end{equation}
La ruta \(\mathbb Z/2\) de
\eqref{p12-83-c20-s8:eq:vnatural-flm-u031} y la ruta \(\mathbb Z/3\) de
\eqref{p12-83-c20-s8:eq:orbifold-triadico-vnatural-u031} son construcciones paralelas
identificadas por teoremas clásicos; ninguna sustituye a la otra.

\subsection{Grupo de automorfismos, carácter y series graduadas}

Una vez construido \(V^\natural\), aparecen dos publicaciones distintas:
\begin{align}
 \operatorname{Aut}(V^\natural)&\cong\mathbb M,
 \label{p12-83-c20-s8:eq:aut-vnatural-monstruo-u031}\\
 \operatorname{ch}_{V^\natural}(\tau)
 &=J(\tau)=j(\tau)-744
 =q^{-1}+196884q+21493760q^2+\cdots.
 \label{p12-83-c20-s8:eq:caracter-vnatural-j-u031}
\end{align}
La primera identifica el grupo de simetrías; la segunda cuenta dimensiones
graduadas. La línea de vacío es
\[
 (V^\natural)_0=\mathbb C\mathbf1,
\]
y la ausencia de raíces de \(\Lambda_{24}\) da
\[
 (V^\natural)_1=0.
\]
En grado dos,
\[
 196884=1+196883=196560+324=54(5\cdot729+1).
\]
Sólo la primera igualdad es una descomposición en dimensiones de
representaciones del Monstruo; las demás son controles aritméticos.

Para \(g\in\mathbb M\), se define la serie de McKay--Thompson
\begin{equation}
 T_g(\tau):=
 \operatorname{Tr}_{V^\natural}
 \left(gq^{L_0-1}\right).
 \label{p12-83-c20-s8:eq:serie-mckay-thompson-u031}
\end{equation}

\begin{theorem}[Moonshine monstruoso]
\label{p12-83-c20-s8:thm:moonshine-monstruoso-u031}
La acción graduada de \(\mathbb M\) sobre \(V^\natural\) produce la familia
\(\{T_g\}_{g\in\mathbb M}\). Los teoremas de Moonshine identifican estas
series con funciones modulares de género cero y prueban sus identidades de
replicabilidad
\cite{ConwayNorton1979,Borcherds1992,FLM1988}.
\end{theorem}

Este teorema tiene como entrada el par
\[
 \bigl(V^\natural,\operatorname{Aut}(V^\natural)\bigr)
\]
y su graduación. No es la consecuencia de una flecha cuyo dominio sea el
grupo abstracto \(\mathbb M\).

Para \(J=T_{1A}\), sean \(P_n\) los polinomios de Faber normalizados por
\[
 P_n(J)=q^{-n}+O(q).
\]
La identidad de Hecke--Faber es
\begin{equation}
 P_n(J)=n\,T_nJ.
 \label{p12-83-c20-s8:eq:hecke-faber-u031}
\end{equation}
Controla todos los grados; el coeficiente 196884 es sólo la primera
manifestación no trivial.

\subsection{Grafo de dependencias del corredor Paley–Witt–Golay–Leech–Moonshine}

\begin{figure}[htbp]
\centering
\begin{tikzcd}[column sep=huge,row sep=large]
 &C_W
   \arrow[dl,"W_{12}"']
   \arrow[d,"N(A_2^{12})"',font=\scriptsize]
   \arrow[dr,phantom,"\substack{\text{ramas}\\\text{distintas}}",
          description,pos=.72,font=\scriptsize]&\\
 M_{12}
 &\Lambda_{24}
   \arrow[d,"V_\Lambda"]
 &(\mathcal G_{24},D,\ell)
   \arrow[d,"W_{24}"]\\
 &V^\natural
   \arrow[dl,"\operatorname{Aut}"']
   \arrow[d,"\operatorname{ch}"]
   \arrow[dr,"g\mapsto T_g"]
 &M_{24}\\
 \mathbb M
 &J
 &\{T_g\}_{g\in\mathbb M}
\end{tikzcd}
\caption{Bifurcación excepcional y tres publicaciones de
\(V^\natural\). La entrada binaria \((\mathcal G_{24},D,\ell)\) es
independiente; no existe una flecha \(W_{24}\to\Lambda_{24}\).}
\label{p12-83-c20-s8:fig:grafo-excepcional-moonshine-u031}
\end{figure}

\begin{criterion}[Criterios de refutación del corredor lunar]
\label{p12-83-c20-s8:crit:refutacion-corredor-lunar-u031}
La construcción queda refutada si ocurre alguno de los hechos siguientes:
\begin{enumerate}
 \item \(\lambda\) no es biyectiva o
       \eqref{p12-83-c20-s8:eq:identidad-entrelazamiento-app-witt-u031} falla;
 \item \(c_*\notin C_W\), \(g_c\) posee puntos fijos o no preserva el
       vecino;
 \item los desplazamientos de
       \eqref{p12-83-c20-s8:eq:desplazamientos-vecino-u031} no pertenecen a \(N_0\);
 \item la elevación reticular \(\widetilde h\) no realiza el automorfismo
       monomial \(h\) en el discriminante;
 \item el peso torcido no es \(4/3\) o la elevación no es de tipo cero;
 \item el orbifold se identifica con \(V^\natural\) sin verificar las
       hipótesis del teorema clásico;
 \item se usa \(W_{24}\) como constructor de Leech;
 \item una igualdad de dimensiones se presenta como construcción de
       \(\mathbb M\);
 \item se serializan indebidamente las publicaciones
       \(\operatorname{Aut}\), \(\operatorname{ch}\) y \(g\mapsto T_g\).
\end{enumerate}
\end{criterion}
